\documentclass[10pt,a4paper]{amsart}
\usepackage[margin=19mm]{geometry}
\usepackage{amsmath,amssymb,mathtools}
\usepackage[scr=rsfs,cal=euler]{mathalfa}
\usepackage{xfrac}
\usepackage[unicode,hidelinks,bookmarks=false]{hyperref}
\newcommand{\ie}{i.e.\ }
\newcommand{\cf}{cf.\ }
\newcommand{\resp}{respectively\ }
\newcommand{\Subsection}{\subsection}
\providecommand{\nobreakdash}{\nobreak\hbox{-}}
\setlength{\emergencystretch}{2em}
\multlinegap0pt
\usepackage{amssymb}
\newtheorem{prop}{Proposition}
\newtheorem{lem}{Lemma}
\title{Dual graphs and generating sequences of non-divisorial valuations on two-dimensional function fields}
\author{Charles Li}
\date{Source: arXiv:1311.1542v3 (CC BY 4.0)}
\begin{document}
\maketitle

\noindent\emph{Source and licence.} Dual graphs and generating sequences of non-divisorial valuations on two-dimensional function fields. Version arXiv:1311.1542v3 (CC BY 4.0), distributed by arXiv under the Creative Commons Attribution 4.0 International licence, http://creativecommons.org/licenses/by/4.0/, as recorded in the arXiv licence metadata for this exact version and shown on the abstract page https://arxiv.org/abs/1311.1542v3. Full licence notice: \texttt{licenses/arxiv-1311.1542-CC-BY-4.0.txt}.\par\smallskip

\noindent\textit{Editorial scope.} Counted unit: the article's lem strict (source0.tex lines 825--831). The article's complete original proof of it is delivered with it (source0.tex lines 833--978, one \texttt{proof} environment of 146 lines). No mathematics is written, repaired or re-derived: the statement and the proof are the article's own lines, in the article's own notation, and no retained line is substituted or re-flowed.  Retained uncounted as the source-local context the proof needs: lines 408--414; lines 444--455; lines 470--491; lines 505--516; lines 543--543; lines 584--586; lines 777--788; lines 810--816; lines 1051--1071.  Nothing was omitted from any retained span, so the package carries no omission record.  No retained line contains a citation, so the wrapper carries no bibliography and no bibliography entry.  No retained line contains a figure environment, an image-inclusion command or drawing code, so no image is delivered and the package declares no asset dependency.  Editorial formatting only, in addition: an AMS \texttt{amsart} preamble that restates the article's own declarations and macros used by the retained lines, namely the environments \texttt{prop}, \texttt{lem} on separate counters, together with the standard packages those lines need.  The retained environments print in this document as Proposition 1, Lemma 1 in order of appearance, whereas the article prints the same items with the numbering of its own full document.  The complete native source of the article as distributed in the arXiv e-print for this exact version is bundled unchanged as \texttt{sources/100097/source0.tex}.
\begin{prop} \label{condet} (Determinant Formula)

Given a continued fraction $\displaystyle [a_1, \ldots, a_n]$, we have:
\begin{equation*}
\lambda_n\mu_{n-1}-\lambda_{n-1}\mu_n = (-1)^{n}
\end{equation*}
\end{prop}

\begin{prop} \label{pq}
Let $\beta_i'=p_i/q_i$. Then,
\begin{equation*}
\left\{
\begin{array}{l}
p_i=\lambda_{m_i} + \lambda_{m_i-1}
\\\\
q_i=\mu_{m_i} + \mu_{m_i-1}
\end{array}
\right.
\end{equation*}
\end{prop}

\begin{lem} \label{XY}
We have the following relationships describing $\left(X,Y\right)$ in terms of $\left(\tilde X,\tilde Y\right)$:
\begin{equation*}
\left\{
\begin{array}{llll}
\text{for } n \text{ even,} & X = \tilde X^{\mu_{n-1}} \tilde Y^{\mu_n}  & \text{and} & Y = \tilde X^{\lambda_{n-1}} \tilde Y^{\lambda_n}
\\\\
\text{for } n \text{ odd,} & X = \tilde X^{\mu_n} \tilde Y^{\mu_{n-1}}  & \text{and} & Y= \tilde X^{\lambda_n} \tilde Y^{\lambda_{n-1}}
\end{array}
\right.
\end{equation*}
More generally,
\begin{equation*}
\left\{
\begin{array}{llll}
\text{for } j \text{ even,} & X = \xi_j^{\mu_{j-1}} \zeta_j^{\mu_j}  & \text{and} & Y = \xi_j^{\lambda_{j-1}} \zeta_j^{\lambda_j}
\\\\
\text{for } j \text{ odd,} & X = \xi_j^{\mu_j} \zeta_j^{\mu_{j-1}}  & \text{and} & Y= \xi_j^{\lambda_j} \zeta_j^{\lambda_{j-1}}
\end{array}
\right.
\end{equation*}
\end{lem}

\begin{lem} \label{XYtilde}
We have the following relationships describing $\left(\tilde X,\tilde Y\right)$ in terms of $\left(X,Y\right)$:
\begin{equation*}
\left\{
\begin{array}{llll}
\text{for } n \text{ even,} & \tilde X = X^{\lambda_n}/Y^{\mu_n} & \text{and} & \tilde Y = Y^{\mu_{n-1}}/X^{\lambda_{n-1}}
\\\\
\text{for } n \text{ odd,} & \tilde X = X^{\lambda_{n-1}}/Y^{\mu_{n-1}} & \text{and} & \tilde Y = Y^{\mu_n}/X^{\lambda_n}
\end{array}
\right.
\end{equation*}
\end{lem}

\section{generating sequences} \label{sectgs}

\begin{equation}  \label{genrecur}
Q_i=Q_{i-1}^{q_{i-1}}-\sum_{h=1}^{\delta_i} u_h^{(i)} \prod_{j=0}^{i-2} Q_j^{\gamma_{j,h}^{(i)}}
\end{equation}

\begin{lem} \label{beta'-1} If $\nu$ is not Type 2, 3 or 4.2, then at $\Sigma(i,0,1)$, where $2\leq i \leq g'$:
\begin{equation*}
\left\{
\begin{array}{l}
\displaystyle \nu(X_i)=\frac{1}{q_1\cdots q_{i-1}}\beta_0
\\\\
\displaystyle \nu(Y_i)=\frac{1}{q_1\cdots q_{i-1}}\left(\beta_i'-1\right)\beta_0
\end{array}
\right.
\end{equation*}
If $\nu$ is Type 2 or 3, then the formulas hold except for $\nu(Y_g)$. If $\nu$ is Type 4.2, then the formulas hold except for $\nu(Y_{g+1})$. See Lemma \ref{genvalue}.
\end{lem}

\begin{lem} \label{leveli+}
For $0\leq i\leq g'-1$, $Q_i$ transforms to:
\begin{equation*}
Q_i=X_{j}^eu
\end{equation*}
for some $e\in\mathbb{N}$, and where $i+1\leq j\leq g'$, and $u$ is a unit.
\end{lem}

\begin{lem} \label{genvalue}
The value of $Q_i$ is:
\begin{equation} \label{usualgenvalue}
\nu(Q_i)=q_{i-1}\nu(Q_{i-1})+\frac{1}{q_1 \cdots q_{i-1}}(\beta_i'-1)\beta_0
\end{equation}
where $2\leq i < g'$. This holds for all $i\geq 2$ in the Type 1 case ($g'=\infty$). This also holds for $i=g'$ in the Type 4.1 case.

For Types 2, 3 and 4.2,
\begin{equation*}
\nu(Q_{g'})=q_{g'-1} \nu(Q_{g'-1})+\nu(Y_{g'})
\end{equation*}
where the jump values are:
\begin{equation*}
\begin{array}{ll}
\text{Type 2:} & \displaystyle \nu(Y_g)= \frac{1}{q_1\cdots q_{g-1}}\left(\tilde \beta_g -1\right) \text{ , where } \tilde \beta_{g}\in\mathbb{R}\setminus\mathbb{Q} \\
\text{Type 3:} & \displaystyle \nu(Y_g)= \left(\lambda_{m-1}^{(g)}-\mu_{m-1}^{(g)},\lambda_{m-2}^{(g)}-\mu_{m-2}^{(g)}\right) \text{ , where } m=m_g \\
\text{Type 4.2:} & \displaystyle \nu(Y_{g+1})=\left(1,\frac{n}{q_1 \cdots q_g}\right) \text{ , where } n\in\mathbb{Z}
\end{array}
\end{equation*}
Also, in the Type 3 case: $\beta_0=q_1\cdots q_{g-1}(\mu_{m-1},\mu_{m-2})$, where $m=m_g$.
\end{lem}

\begin{lem} \label{strict}
For $2\leq i\leq g'$, $Q_i$ transforms to:
\begin{equation*}
Q_i=\tilde X_{i-1}^{f_1}\tilde Y_{i-1}^{f_2}\left(\tilde Y_{i-1}-c_i\tilde X_{i-1}\right)u
\end{equation*}
for some $f_1,f_2\in\mathbb{N}$, and where $u$ is a unit and $c_i\in k$ is the residue of $\tilde Y_{i-1}/\tilde X_{i-1}$.
\end{lem}

\begin{proof}
First, we show $\displaystyle \nu(Q_1)=\frac{p_1}{q_1}\beta_0$. Let $m=m_1$. By Lemma \ref{XYtilde} and the fact that $\nu(\tilde X_1)=\nu(\tilde Y_1)$, we get:
\begin{equation*}
\nu\left(\frac{x^{\lambda_m}}{y^{\mu_m}}\right)=\nu\left(\frac{y^{\mu_{m-1}}}{x^{\lambda_{m-1}}}\right)
\end{equation*}
or
\begin{equation*}
\nu\left(\frac{x^{\lambda_{m-1}}}{y^{\mu_{m-1}}}\right)=\nu\left(\frac{y^{\mu_{m}}}{x^{\lambda_{m}}}\right)
\end{equation*}
depending on whether $m$ is even or odd, respectively. In either case,
\begin{equation*}
\left(\lambda_m+\lambda_{m-1}\right)\nu(x)=\left(\mu_m+\mu_{m-1}\right)\nu(y)
\end{equation*}
By Proposition \ref{pq}, $\lambda_m+\lambda_{m-1}=p_1$ and $\mu_m+\mu_{m-1}=q_1$, so $\displaystyle \nu(Q_1)=\frac{p_1}{q_1}\beta_0$.

Now, use induction on $i$ to prove the lemma. By Equation (\ref{genrecur}) in Section \ref{sectgs},
\begin{equation*}
Q_2=Q_1^{q_1}-u_1Q_0^{\gamma_{0,1}}
\end{equation*}
and $\displaystyle q_1\nu(Q_1) = \gamma_{0,1} \beta_0$. Thus, $\gamma_{0,1}=p_1$ and $\displaystyle Q_2 = y^{q_1} - u_1x^{p_1}$.

Assume $m=m_1$ is even for concreteness and apply Lemma \ref{XY}:
\begin{equation*}
Q_2 = \tilde X_1^{q_1\lambda_{m-1}} \tilde Y_1^{q_1\lambda_{m}} - u_1 \tilde X_1^{p_1\mu_{m-1}} \tilde Y_1^{p_1\mu_{m}}
\end{equation*}
Substituting $q_1 = \mu_m + \mu_{m-1}$ and $p_1=\lambda_m + \lambda_{m-1}$, then factoring, we get:
\begin{equation*}
Q_2 = \tilde X_1^{f_1} \tilde Y_1^{f_2}\left(\tilde Y_1^{\lambda_m\mu_{m-1}-\lambda_{m-1}\mu_m} -u_1 \tilde X_1^{\lambda_m\mu_{m-1}-\lambda_{m-1}\mu_m} \right)
\end{equation*}
where $f_1=\lambda_{m-1}\mu_m+\lambda_{m-1}\mu_{m-1}$ and $f_2=\lambda_{m-1}\mu_m+\lambda_m\mu_m$.
Now apply Proposition \ref{condet} to simplify the inside of the parentheses.
\begin{equation*}
Q_2 = \tilde X_1^{f_1} \tilde Y_1^{f_2}\left(\tilde Y_1 -u_1 \tilde X_1 \right)
\end{equation*}

This shows the base step for the induction in the even $m_1$ case. The odd $m_1$ case is similar and the details are omitted. Now we show $Q_{i+1}$ behaves nicely given the inductive hypothesis up to level $i$. The plan of attack is to compute the total transforms of the $\{Q_j\}_{j=0}^{i}$ at $\Sigma(i,0,1)$, then use these calculations to prove the conclusion for $Q_{i+1}$.

For $Q_0=x$ and $Q_1=y$, applying Lemma \ref{XY} followed by a $z$-blowup and then repeating the process shows that $Q_0=X_j^{f_{0,j}}u_{0,j}$ and $Q_1=X_j^{f_{1,j}}u_{1,j}$ at $\Sigma(j,0,1)$ for $2\leq j \leq g'$, where $u_{0,j}$ and $u_{1,j}$ are units, and $f_{0,j},f_{1,j}\in\mathbb{N}$.

For $2\leq j \leq i$, at level $\Sigma(j,0,0)$, assume:
\begin{equation*}
Q_j=\tilde X_{j-1}^{f_1} \tilde Y_{j-1}^{f_2} \left(\tilde Y_{j-1} - c_j \tilde X_{j-1}\right) u
\end{equation*}
where the positive integers $f_1$ and $f_2$ are understood to be different for each $j$, but we suppress any subscripts to indicate as such for the sake of lucidity. Although the units $u$ might vary with each blowup, they will remain units, so we suppress notation here as well.

Performing the next $z$-blowup to get to level $\Sigma(j,0,1)$:

\begin{equation*}
\begin{array}{rl}
\displaystyle Q_j= & X_{j}^{f_1}X_{j}^{f_2}\left(Y_{j}+c_j\right)^{f_2}X_{j}Y_{j} u
\\\\
\displaystyle = & X_{j}^{f_1+f_2+1}Y_{j} u
\end{array}
\end{equation*}
where the $\left(Y_{j}+c_j\right)^{f_2}$ was absorbed into the unit. Let $f_3=f_1+f_2+1$ for simplicity.

Let $m=m_j$ and assume $m_j$ is odd for concreteness. The even case is similar. Tracing blowups to levels $\Sigma(j+1,0,0)$ and $\Sigma(j+1,0,1)$, we have:
\begin{equation*}
\begin{array}{rl}
\displaystyle Q_j=&\tilde X_{j}^{f_3\mu_{m}} \tilde Y_{j}^{f_3\mu_{m-1}} \tilde X_{j}^{\lambda_{m}} \tilde Y_{j}^{\lambda_{m-1}}u
\\\\
\displaystyle =&\tilde X_{j}^{f_3\mu_{m}+\lambda_m} \tilde Y_{j}^{f_3\mu_{m-1}+\lambda_{m-1}}u
\\\\
\displaystyle =& X_{j+1}^{f_3(\mu_{m}+\mu_{m-1})+\lambda_m+\lambda_{m-1}} (Y_{j+1}+c_{j+1})^{f_3\mu_{m-1}+\lambda_{m-1}}u
\\\\
\displaystyle =& X_{j+1}^{f_4}u
\end{array}
\end{equation*}
where the $(Y_{j+1}+c_{j+1})^{f_3\mu_{m-1}+\lambda_{m-1}}$ was absorbed into $u$ and the exponent of $X_{j+1}$ was relabeled as $f_4$.

Tracing the blowups further, it is easy to see that $Q_j$ will have a total transform of the form $\displaystyle X_{k}^{e_k}u$ at level $\Sigma(k,0,1)$ for all $j+1\leq k \leq g'$, where $e_k\in\mathbb{N}$ and $u$ is a unit. This proves Lemma \ref{leveli+} once the proof of Lemma \ref{strict} is complete.

Now, we return to $\displaystyle Q_{i+1}=Q_i^{q_i}-\sum_{h=1}^{\delta_{i+1}} u_h \prod_{j=0}^{i-1}Q_j^{\gamma_{j,h}}$. By the previous discussion, at level $\Sigma(i,0,1)$, $Q_i=X_{i}^{f_1}Y_{i}u$ for some $f_1\in\mathbb{N}$, and so $Q_i^{q_i}=X_i^{q_if_1}Y_i^{q_i}u$, ignoring changes to the unit. Also by the previous discussion, each of the $\prod Q_j^{\gamma_{j,h}}$ will transform to $X_{i}^{f_2}u_h$ for all $h$, with the same power $f_2$ and differing only in the units $u_h$; just transform each $Q_j^{\gamma_{j,h}}$, then collect the $X_{i}$ factors together. The common value $\nu(\prod Q_j^{\gamma_{j,h}})=q_i\nu(Q_i)$ ensures the same $f_2$ for all $h$. Now, factor out the $X_{i}^{f_2}$ in the total transform of $\sum u_h \prod Q_j^{\gamma_{j,h}}$ and set $u'=\sum_{h=1}^{\delta_{i+1}} u_h$. We get:
\begin{equation} \label{qif1}
Q_{i+1}= X_{i}^{q_if_1}Y_{i}^{q_i}u - X_{i}^{f_2}u'
\end{equation}

Notice that $u'\neq 0$ or else $Q_{i+1}$ would not give a jump in value. Note that the two terms on the right both have value $q_i \nu(Q_i)$. Using Lemma \ref{beta'-1} and taking values we have:
\begin{equation*}
\nu\left(X_{i}^{q_if_1}Y_{i}^{q_i}u\right)=q_if_1 \cdot \frac{1}{q_1\cdots q_{i-1}}\beta_0 + q_i\cdot \frac{1}{q_1\cdots q_{i-1}}(\beta_i'-1)\beta_0
\end{equation*}
\begin{equation*}
\nu(X_i^{f_2}u')=f_2\cdot \frac{1}{q_1 \cdots q_{i-1}}\beta_0
\end{equation*}
Set equal and clear the $\prod q_j$ and $\beta_0$. Since $\beta_i'=p_i/q_i$, we get:
\begin{equation*}
q_i f_1=f_2-p_i+q_i
\end{equation*}
Substituting for $q_if_1$ in (\ref{qif1}), we have:
\begin{equation*}
\begin{array}{rl}
\displaystyle Q_{i+1}= & X_{i}^{f_2-p_i+q_i}Y_{i}^{q_i}u - X_{i}^{f_2}u'
\\\\
= & X_{i}^{f_2-p_i+q_i}\left[Y_{i}^{q_i} - c_{i+1}X_{i}^{p_i-q_i}\right]u
\end{array}
\end{equation*}
where we factored out $u$ and set $c_{i+1}\cong u'/u$, where $c_{i+1}\in k$. This is possible because the ring $R_{\Sigma(i,0,1)}$ is localized at its maximal ideal and the residue field is isomorphic to $k$.

Notice it is only necessary to show that $Y_{i}^{q_i} - c_{i+1}X_{i}^{p_i-q_i}$ transforms to the form: $\tilde X_{i}^{f_3} \tilde Y_{i}^{f_4} \left(\tilde Y_{i}-c_{i+1} \tilde X_{i}\right)$, where $f_3,f_4\in\mathbb{N}$. Let $m=m_i$ and assume it is even for concreteness. The odd case is similar. Apply Lemma \ref{XY} using $X=X_i$, $Y=Y_i$, $\beta=\left[a_1^{(i)}-1,a_2^{(i)},\cdots,a_{m}^{(i)}\right]$. The subtraction of 1 in $a_1^{(i)}-1$ represents one less blowup since the $(X_i,Y_i)$ parameters occur after a $z$-blowup, which we have to account for. Let $\lambda_j/\mu_j$ be the convergents of $\beta_i'$. Then the convergents of $\beta$ will be $\lambda_j/\mu_j-1=(\lambda_j-\mu_j)/\mu_j$. Applying Lemma \ref{XY} gives:
\begin{equation*}
Y_{i}^{q_i} - c_{i+1}X_{i}^{p_i-q_i} = \tilde X_i^{q_i(\lambda_{m-1}-\mu_{m-1})} \tilde Y_i^{q_i(\lambda_m-\mu_m)} - c_{i+1} \tilde X_i^{(p_i-q_i)\mu_{m-1}} \tilde Y_i^{(p_i-q_i)\mu_m}
\end{equation*}
Expanding the exponents, we get:
\begin{equation*}
\begin{array}{rl}
q_i(\lambda_{m-1}-\mu_{m-1}) & = (\mu_{m-1}+\mu_m)(\lambda_{m-1}-\mu_{m-1})
\\\\
 & =\lambda_{m-1}\mu_{m-1}+\lambda_{m-1}\mu_m-\mu_{m-1}^2-\mu_{m-1}\mu_m
\end{array}
\end{equation*}
\begin{equation*}
\begin{array}{rl}
q_i(\lambda_m-\mu_m) & = (\mu_{m-1}+\mu_m)(\lambda_m-\mu_m)
\\\\
& =\lambda_m\mu_{m-1}+\lambda_m\mu_m-\mu_{m-1}\mu_m-\mu_m^2
\end{array}
\end{equation*}
\begin{equation*}
\begin{array}{rl}
(p_i-q_i)\mu_{m-1} & = (\lambda_{m-1}+\lambda_m-\mu_{m-1}-\mu_m)\mu_{m-1}
\\\\
& =\lambda_{m-1}\mu_{m-1}+\lambda_m\mu_{m-1}-\mu_{m-1}^2-\mu_{m-1}\mu_m
\end{array}
\end{equation*}
\begin{equation*}
\begin{array}{rl}
(p_i-q_i)\mu_m & =(\lambda_{m-1}+\lambda_m-\mu_{m-1}-\mu_m)\mu_m
\\\\
& =\lambda_{m-1}\mu_m+\lambda_m\mu_m-\mu_{m-1}\mu_m-\mu_m^2
\end{array}
\end{equation*}

Let:
\begin{equation*}
f_3=\lambda_{m-1}\mu_{m-1}+\lambda_{m-1}\mu_m-\mu_{m-1}^2-\mu_{m-1}\mu_m
\end{equation*}
and
\begin{equation*}
f_4=\lambda_{m-1}\mu_m+\lambda_m\mu_m-\mu_{m-1}\mu_m-\mu_m^2
\end{equation*}
Factor out $\tilde X^{f_3}\tilde Y^{f_4}$:
\begin{equation*}
Y_{i}^{q_i} - c_{i+1}X_{i}^{p_i-q_i} = \tilde X_i^{f_3} \tilde Y_i^{f_4}\left(\tilde Y_i^{\lambda_m \mu_{m-1} - \lambda_{m-1} \mu_m} - c_{i+1}\tilde X_i^{\lambda_m\mu_{m-1}-\lambda_{m-1}\mu_m}\right)
\end{equation*}
By Proposition \ref{condet}, $\lambda_m\mu_{m-1}-\lambda_{m-1}\mu_m=1$ and we are done.
\end{proof}

\end{document}
